%%%%%%%%%%  scan idx 577  |  folio 564  %%%%%%%%%%
\SGAPrintPage{564}%

where the vertical arrows are the canonical homomorphisms \(K^{\bullet}\to K_\bullet\) (IV 2.5.1).

\sgaStableTarget{sga6.e10.sec.7.7}\subsection*{7.7. Specialization for \(K^\bullet\) and \(\Gr^\bullet\).}
When \(X_t\) and \(X_s\) are regular (for example, when \(f\) is smooth), the two preceding vertical arrows are isomorphisms (IV 2.5), and the specialization homomorphism can then be interpreted as a homomorphism
\sgaStableTarget{sga6.e10.eq.7-7-1}\begin{equation}\tag{7.7.1}
  \sigma^\bullet:K^\bullet(X_t)\longrightarrow K^\bullet(X_s),
\end{equation}
which can also be described as the unique homomorphism making the following diagram commutative
\sgaStableTarget{sga6.e10.eq.7-7-2}\begin{equation}\tag{7.7.2}
\begin{tikzcd}[column sep=large,row sep=large]
K^\bullet(X) \arrow[r,"j^*"] \arrow[d,"i^*"'] & K^\bullet(X_t) \arrow[dl,"\sigma^\bullet"]\\
K^\bullet(X_s) & {.}
\end{tikzcd}
\end{equation}
Since, in this last diagram, \(i^*\) and \(j^*\) are ring homomorphisms (IV 2.7 b), with \(j^*\) surjective, it follows that \emph{the specialization homomorphism} (7.7.1) \emph{is a ring homomorphism}. For the same reason, \emph{if} \(X\) \emph{is separated}, \emph{the homomorphism} (7.7.1) \emph{is even a homomorphism of} \(\lambda\)-\emph{rings}. To make this precise, one need only note that a regular separated noetherian scheme is divisorial (II 2.2.7.1), so its \(K^\bullet\) is identified with its ``naive'' \(K^\bullet\) (IV 2.9 (ii)), which is naturally endowed with a \(\lambda\)-ring structure (VI 3.2). On the other hand, the regularity of \(X_t\) and \(X_s\) implies that of \(X\) (EGA \(0_{\mathrm{IV}}\) 17.1.8), and the homomorphisms \(i^*:K^\bullet(X)\to K^\bullet(X_s)\) and \(j^*:K^\bullet(X)\to K^\bullet(X_t)\) are \(\lambda\)-homomorphisms, so the same is true of (7.7.1), which is obtained from \(i^*\) by passage to the quotient.

%%%%%%%%%%  scan idx 578  |  folio 565  %%%%%%%%%%
\SGAPrintPage{565}%

Note also that, since \(X\) is regular and hence normal, its connected components coincide with its irreducible components, which dominate \(S\), so that the homomorphism \(H^0(X,\mathbb Z)\to H^0(X_t,Z)\) induced by the inclusion \(X_t\to X\) is bijective. Consequently the homomorphism \(H^0(X,\mathbb Z)\to H^0(X_s,\mathbb Z)\) induced by the inclusion \(X_s\to X\) can also be interpreted as a homomorphism
\sgaStableTarget{sga6.e10.eq.7-7-3}\begin{equation}\tag{7.7.3}
  H^0(X_t,\mathbb Z)\longrightarrow H^0(X_s,\mathbb Z),
\end{equation}
which again deserves the name ``specialization homomorphism''. One then checks immediately from the definitions that the homomorphism (7.7.1) is compatible with the \emph{augmented algebra} structures over \(H^0(X_t,\mathbb Z)\) and \(H^0(X_s,\mathbb Z)\) respectively on the two sides, defined in (IV 2.8), taking account of the preceding homomorphism (7.7.3).

Still supposing that \(X\) is separated, the preceding discussion implies in particular that (7.7.1) induces a homomorphism of graded rings
\sgaStableTarget{sga6.e10.eq.7-7-4}\begin{equation}\tag{7.7.4}
  \sigma^\bullet:\Gr^\bullet(X_t)\longrightarrow\Gr^\bullet(X_s),
\end{equation}
and, using the notion of Chern class of (V 6), one obtains a commutative diagram
\[
\begin{tikzcd}[column sep=large,row sep=large]
K^\bullet(X_t) \arrow[r,"c_{X_t}"] \arrow[d] &
\Ch\,\Gr^\bullet(X_t) \arrow[d]\\
K^\bullet(X_s) \arrow[r,"c_{X_s}"] &
\Ch\,\Gr^\bullet(X_s).
\end{tikzcd}
\]

\sgaStableTarget{sga6.e10.sec.7.8}\subsection*{7.8. Specialization for \(\Pic\).}
Retain the hypotheses of 7.7: \(X\) separated, \(X_t\) and \(X_s\) regular. Taking account of the canonical isomorphism (5.3.2):
\[
  \Gr^1\simeq \Pic ,
\]
the homomorphism (7.7.4) therefore gives a canonical homomorphism
\sgaStableTarget{sga6.e10.eq.7-8-1}\begin{equation}\tag{7.8.1}
  \Pic(X_t)\longrightarrow \Pic(X_s).
\end{equation}

%%%%%%%%%%  scan idx 579  |  folio 566  %%%%%%%%%%
\SGAPrintPage{566}%

An equivalent description of this last homomorphism is as follows: the hypothesis that \(X\) is regular implies that these local rings are factorial (by the Auslander--Buchsbaum theorem, EGA IV 21.11.1), and hence that every invertible sheaf \(\underline{L}_t\) on \(X_t\) extends to an invertible sheaf \(\underline{L}\) on \(X\) (EGA IV 21.6.11). One then checks at once that the image under (7.8.1) of \(\cl(\underline{L}_t)\in\Pic(X_t)\) is \(\cl(\underline{L}_s)\in\Pic(X_s)\), where, of course, one puts \(\underline{L}_s=\underline{L}|X_s\).

In fact, using only the hypothesis that the local rings of \(X\) are factorial and that \(X_s\) is locally integral (without hypotheses of regularity or separatedness), one shows that \(\cl(\underline{L}_s)\) does not depend on the chosen extension \(\underline{L}\) of \(\underline{L}_t\) to an invertible Module on \(X\), so that one obtains a homomorphism (7.8.1) directly in this way. To verify the asserted independence, one is reduced to proving that, if \(\underline{L}_t\simeq\underline{O}_{X_t}\), then \(\underline{L}_s\simeq\underline{O}_{X_s}\). Now the hypothesis means that \(\underline{L}\simeq\underline{O}_X(D)\), with \(D\) a divisor on \(X\) having support \(\subset X_s\) (cf. EGA IV 21.2.11), and one reduces at once to the case where \(X_s\) is \emph{integral}, hence \(D=nX_s\); but then \(D=\operatorname{div}(g^n)\) (\(g\) a uniformizer of \(A\)), so \(\underline{O}_X(D)\simeq\underline{O}_X\) and a fortiori \(\underline{L}_s\simeq\underline{O}_{X_s}\).

\sgaStableTarget{sga6.e10.sec.7.9}\subsection*{7.9. Specialization and numerical equivalence.}
Suppose that \(f\) is proper, and \(X_s\) and \(X_t\) are regular, and consider the two specialization maps (7.4.1) and (7.7.4):
\sgaStableTarget{sga6.e10.eq.7-9-1}\begin{equation}\tag{7.9.1}
\begin{gathered}
  \sigma_\bullet:\Gr_\bullet(X_t)\longrightarrow\Gr_\bullet(X_s)\\
  \sigma^\bullet:\Gr^\bullet(X_t)\longrightarrow\Gr^\bullet(X_s).
\end{gathered}
\end{equation}
Since the homomorphism
\sgaStableTarget{sga6.e10.eq.7-9-2}\begin{equation}\tag{7.9.2}
  \sigma:K^\bullet(X_t)\longrightarrow K^\bullet(X_s)
\end{equation}

%%%%%%%%%%  scan idx 580  |  folio 567  %%%%%%%%%%
\SGAPrintPage{567}%

from which they are obtained is compatible with the ring structures, using (7.5.1) one finds that
\sgaStableTarget{sga6.e10.eq.7-9-3}\[
  \langle\sigma_\bullet(x_t),\sigma^\bullet(y_t)\rangle
  =
  \langle x_t,y_t\rangle
  \qquad\text{for }
  x_t\in\Gr_\bullet(X_t)\quad y_t\in\Gr^\bullet(X_t),
\tag{7.9.3}
\]
with the pairings \(\langle x,y\rangle\) of 4.1 between \(\Gr_\bullet\) and \(\Gr^\bullet\). It follows formally that \emph{if} \(x\) \emph{is an element of} \(\Gr_\bullet(X_t)\) (\emph{resp}. \emph{of} \(\Gr^\bullet(X_t)\)) \emph{such that} \(\sigma_\bullet(x)\) (\emph{resp}. \(\sigma^\bullet(x)\)) \emph{is numerically equivalent to zero} (4.5.1), \emph{then the same is true of} \(x\). The converse is true when (7.9.2) (or, equivalently, either of the homomorphisms (7.9.1) obtained from it by passage to the associated graded groups) is an epimorphism mod.torsion, - a circumstance that occurs only rather exceptionally.

\sgaStableTarget{sga6.e10.para.7-9-4}\textbf{7.9.4.} It is plausible, at least when \(f\) is proper and \emph{smooth}, that the converse in question always holds, i.e., that \(\sigma_\bullet(x)\) (resp. \(\sigma^\bullet(x)\)) is numerically equivalent to zero as soon as the same is true of \(x\). This would imply that the homomorphisms (7.9.1) induce \emph{monomorphisms} on the corresponding quotient groups by numerical equivalence, \(\Gr^{\mathrm{num}}_\bullet\) and \(\Gr_{\mathrm{num}}^\bullet\). At least the preceding discussion implies that the group \(\Gr_\bullet^{\mathrm{num}}(X_t)\) is isomorphic to a quotient of a subgroup of \(\Gr_\bullet^{\mathrm{num}}(X_s)\) (and in particular has rank at most that of \(\Gr_\bullet^{\mathrm{num}}(X_s)\)), and likewise for \(\Gr_{\mathrm{num}}^\bullet\), and also for \(K^\bullet_{\mathrm{num}}\) itself.

\sgaStableTarget{sga6.e10.sec.7.10}\subsection*{7.10. Compatibility with direct and inverse images}
Let
\[
  g:X\longrightarrow Y
\]
be a morphism of noetherian schemes flat over \(S\). Denote by \(\sigma_X\) (resp. \(\sigma_Y\)) the specialization homomorphism for \(X\) (resp. \(Y\)), and by \(i_X,j_X\) (resp. \(i_Y,j_Y\))

%%%%%%%%%%  scan idx 581  |  folio 568  %%%%%%%%%%
\SGAPrintPage{568}%

the morphisms \(i,j\) for \(X\) (resp. \(Y\)) from the special and generic fibers to the ambient scheme. Denote by
\[
  g_t:X_t\to Y_t,\qquad g_s:X_s\to Y_s
\]
the induced morphisms, which therefore give rise to the Cartesian squares
\sgaStableTarget{sga6.e10.eq.7-10-0}\begin{equation}\tag{7.10.0}
\begin{tikzcd}[column sep=large,row sep=large]
X_t \arrow[r,"g_t"] \arrow[d,"j_X"'] & Y_t \arrow[d,"j_Y"'] &
X_s \arrow[r,"g_s"] \arrow[d,"i_X"'] & Y_s \arrow[d,"i_Y"']\\
X \arrow[r,"g"] & Y &
X \arrow[r,"g"] & Y .
\end{tikzcd}
\end{equation}

First suppose that \(g\) is \emph{proper}. I claim that the following diagram then commutes:
\sgaStableTarget{sga6.e10.eq.7-10-1}\begin{equation}\tag{7.10.1}
\begin{tikzcd}[column sep=large,row sep=large]
K_\bullet(X_t) \arrow[r,"g_{t*}"] \arrow[d,"\sigma_X"'] &
K_\bullet(Y_t) \arrow[d,"\sigma_Y"]\\
K_\bullet(X_s) \arrow[r,"g_{s*}"] &
K_\bullet(Y_s)
\end{tikzcd}
\end{equation}
This result generalizes 7.5 (since, for \(X=S\), the arrows of (7.3.2) are identified with the identity map of \(\mathbb Z\)). Note that, when \(X\) and \(Y\) satisfy the conditions of 7.7, i.e., have regular fibers, \(g_t\) and \(g_s\) automatically have finite Tor-dimension, and the diagram can be interpreted as a commutative diagram with \(K^\bullet\), which we leave the reader to rewrite.

To prove the commutativity of (7.10.1), one sees at once, using the definitions of \(\sigma_X\) and \(\sigma_Y\), that one is reduced to proving the commutativity of the two front squares of the prism
\sgaStableTarget{sga6.e10.eq.unnamed}\[
\begin{tikzcd}[column sep=large,row sep=large]
K_\bullet(X_t) \arrow[rr,"g_{t*}"] \arrow[dd] && K_\bullet(Y_t) \arrow[dd]\\
& K_\bullet(X) \arrow[ul,"j_X^*" description] \arrow[dl,"i_X^*" description] \arrow[rr,"g_*"] &&
K_\bullet(Y) \arrow[ul,"j_Y^*" description] \arrow[dl,"i_Y^*" description]\\
K_\bullet(X_s) \arrow[rr,"g_{s*}"] && K_\bullet(Y_s).
\end{tikzcd}
\tag{*}
\]
